\section{Вычислительные методы линейной алгебры}
\label{sec:numerical-linear-algebra}

\subsection{Задачи вычислительной линейной алгебры}

Основные задачи:
\begin{itemize}
    \item решение систем линейных алгебраических уравнений $Ax=b$;
    \item вычисление собственных значений и собственных векторов;
    \item задачи наименьших квадратов;
    \item матричные разложения и оценка обусловленности.
\end{itemize}

Численный метод должен учитывать округление, накопление ошибок, структуру и
разреженность матрицы. Ошибка решения определяется не только алгоритмом, но и
обусловленностью самой задачи.

\subsection{Метод Гаусса}

Прямой ход метода Гаусса последовательными элементарными преобразованиями строк
приводит систему к верхнетреугольному виду
\[
    Ux=\widetilde b.
\]
Затем выполняется обратная подстановка:
\[
    x_i=\frac{1}{u_{ii}}
    \left(\widetilde b_i-\sum_{j=i+1}^{n}u_{ij}x_j\right).
\]

Число арифметических операций имеет порядок $O(n^3)$, память для плотной матрицы
--- $O(n^2)$. Без выбора главного элемента метод может быть численно неустойчив.
Обычно применяют частичный выбор: на $k$-м шаге выбирают максимальный по модулю
элемент в столбце среди строк $k,\ldots,n$ и меняют строки местами.

\subsection[LU-разложение]{$LU$-разложение}

Метод Гаусса эквивалентен факторизации
\[
    PA=LU,
\]
где $P$ --- матрица перестановки, $L$ --- нижнетреугольная матрица с единичной
диагональю, $U$ --- верхнетреугольная. Решение выполняется в два этапа:
\[
    Ly=Pb,
    \qquad
    Ux=y.
\]

Факторизация стоит $O(n^3)$, а решение для каждого нового правого столбца ---
$O(n^2)$. Поэтому $LU$ особенно выгодно, когда матрица одна, а правых частей
несколько.

Определитель можно вычислить как
\[
    \det A=\det(P)\prod_{i=1}^{n}u_{ii}.
\]

\subsection{Разложение Холецкого}

Если $A$ симметрична и положительно определена, существует единственное
разложение
\[
    A=LL^{T}
\]
с положительной диагональю $L$. Оно примерно вдвое дешевле общего $LU$ и не
требует выбора главного элемента. Формулы:
\[
 l_{ii}=\sqrt{a_{ii}-\sum_{k=1}^{i-1}l_{ik}^2},
\]
\[
 l_{ji}=\frac{1}{l_{ii}}
 \left(a_{ji}-\sum_{k=1}^{i-1}l_{jk}l_{ik}\right),
 \qquad j>i.
\]

Разложение Холецкого широко применяется для нормальных уравнений, ковариационных
матриц и конечно-элементных систем.

\subsection{Метод прогонки}

Для трёхдиагональной системы
\[
 a_i x_{i-1}+b_i x_i+c_i x_{i+1}=d_i
\]
метод прогонки, или алгоритм Томаса, выполняет исключение за $O(n)$ операций и
использует $O(n)$ памяти. Представим
\[
    x_i=\alpha_i x_{i+1}+\beta_i.
\]
Коэффициенты вычисляются прямым ходом, после чего выполняется обратная
подстановка. Устойчивость гарантируется, например, при строгом диагональном
преобладании.

\subsection[QR-разложение]{$QR$-разложение}

Матрица представляется как
\[
    A=QR,
\]
где $Q^TQ=I$, а $R$ верхнетреугольна. Разложение строят преобразованиями
Хаусхолдера или вращениями Гивенса. Оно устойчивее нормальных уравнений в задаче
наименьших квадратов:
\[
    \min_x\|Ax-b\|_2
    \quad\Longrightarrow\quad
    Rx=Q^Tb.
\]
Преобразования Гивенса удобны для разреженных матриц и последовательного
обнуления отдельных элементов.

\subsection{Стационарные итерационные методы}

Пусть $A=D+L+U$, где $D$ --- диагональная часть, $L$ и $U$ --- нижняя и верхняя
строго треугольные части.

Метод Якоби:
\[
    x^{(k+1)}=-D^{-1}(L+U)x^{(k)}+D^{-1}b.
\]
Метод Гаусса--Зейделя:
\[
    (D+L)x^{(k+1)}=b-Ux^{(k)}.
\]
Метод верхней релаксации:
\[
    (D+\omega L)x^{(k+1)}
    =\omega b-\bigl(\omega U+(\omega-1)D\bigr)x^{(k)}.
\]

Для итерации $x^{(k+1)}=Bx^{(k)}+c$ необходимым и достаточным условием сходимости
при любом начальном приближении является
\[
    \rho(B)<1,
\]
где $\rho(B)$ --- спектральный радиус. Достаточное условие для Якоби и Зейделя
--- строгое диагональное преобладание; для симметричной положительно определённой
матрицы метод Зейделя сходится.

\subsection{Методы подпространств Крылова}

Для больших разреженных систем используют пространства
\[
    \mathcal K_k(A,r_0)=
    \operatorname{span}\{r_0,Ar_0,\ldots,A^{k-1}r_0\}.
\]

Метод сопряжённых градиентов применяется к симметричным положительно определённым
матрицам. В точной арифметике он даёт решение не более чем за $n$ шагов, а на
практике скорость зависит от распределения собственных значений. Один шаг
требует умножения матрицы на вектор.

Для несимметричных матриц применяют GMRES, BiCGStab и другие методы. Ключевой
элемент --- предобуславливание:
\[
    M^{-1}Ax=M^{-1}b,
\]
где $M$ приближает $A$, но легко обращается.

\subsection{Собственные значения}

Степенной метод:
\[
    y^{(k+1)}=Ax^{(k)},
    \qquad
    x^{(k+1)}=\frac{y^{(k+1)}}{\|y^{(k+1)}\|}
\]
сходится к собственному вектору, соответствующему доминирующему по модулю
собственному значению, если оно единственно и начальный вектор имеет ненулевую
проекцию на его собственный вектор.

Обратная итерация
\[
    (A-\mu I)y^{(k+1)}=x^{(k)}
\]
находит собственное значение, ближайшее к сдвигу $\mu$. Отношение Рэлея
\[
    \rho(x)=\frac{x^TAx}{x^Tx}
\]
даёт приближение собственного значения.

Для симметричных матриц применяют метод вращений Якоби или $QR$-алгоритм.

\subsection{Обусловленность и невязка}

Число обусловленности в согласованной норме:
\[
    \kappa(A)=\|A\|\,\|A^{-1}\|.
\]
При малом возмущении правой части
\[
    \frac{\|\delta x\|}{\|x\|}
    \lesssim
    \kappa(A)\frac{\|\delta b\|}{\|b\|}.
\]
Малая невязка
\[
    r=b-A\widetilde x
\]
не всегда означает малую ошибку решения: при большой $\kappa(A)$ задача сама
чувствительна.

\subsection{Что важно сказать на экзамене}

Следует разделить прямые и итерационные методы. Для плотных систем стандартны
Гаусс и $LU$, для симметричных положительно определённых --- Холецкий, для
трёхдиагональных --- прогонка, для больших разреженных --- методы Крылова с
предобуславливанием. Обязательно упомянуть устойчивость и обусловленность.

\newpage